\documentclass[a4paper,12pt]{article}
\usepackage[margin=1in]{geometry}
\usepackage{graphicx}
\usepackage{hyperref}

\title{Quantum Error Correction Decoder Benchmark}
\author{Andrew Nestor}
\date{\today}

\begin{document}

\maketitle

\begin{abstract}
Investigating QEC
\end{abstract}

\section{Introduction}

This project focused on building a QEC system, starting with a state based simulation of the 3 qubit repitition code

\section{Background}

This project focused on stabilizer codes.

\subsection{Quantum Error Correcting Codes}
\subsection{Stabilizer Codes}
\subsection{Repitition Codes}
\subsection{Topological Codes}
\subsection{Decoding}

\section{State Simulation}

I started by creating a circuit from scratch.
This was done by simulatingthe full state as a complex vectors, and gates as unitary matrices.

The gates were implemented as part of the \verb|Circuit| class.
I implemented \textbf{H}, \textbf{X}, \textbf{Z}, \textbf{S}, and \textbf{CNOT} gates
All single qubit gates were implemented using \verb|np.kron(G1,G2)| where kron denotes the Kroneker product
For example, to apply \textbf{H} to qubuts 2 and 3 of a 3-qubit state, we construct

\begin{equation}
    H = I \oplus H_2 \oplus H_3 
\end{equation}

Began by focusing on the 3 qubit repitition code. This includes two parity check stabilizers, defined as

\begin{equation}
    S_1 = Z_1 Z_2 \;\;\;\; S_2 = Z_2 Z_3.
\end{equation}

Each of these were measures using a pair of \textbf{CNOT} gates from the target qubits onto an ancilla

\subsection{Measurements}
\subsection{Errors}
\subsection{Decoding}
\subsubsection{Matching}
\subsection{Measuring Logical Errors}
\subsection{Running}

\begin{verbatim}
    Instantiate Gates to create initial code (inc. Stabilizers)
    for shot in shots:    
        for round in rounds:
            add errors
            run gates
            measure syndrome
            reset stabilizers
        decode syndrome
        predict errors

    get logical error
        
\end{verbatim}

\subsection{Validation}


\subsection{Experiments}

\begin{itemize}
    \item Behaviour as $p_m \rightarrow 0$, $p_x \rightarrow 0$
    \item Different noise models (weights, types, etc)
    \item Different rounds vs probabilities
    \item sim speed
\end{itemize}


\section{Stabilizer Simulation}

\section{Noise Models}

\section{Decoding}
\subsection{Matching - Blossom Algorithm}
\subsection{Optimisation}
\subsection{Heuristic Methods}
\subsection{Machine Learning}

\section{Results}

\section{Conclusions and Future Work}

\subsection{Conclusions}
\subsection{Future Work}
\begin{itemize}
    \item Adapting the NetworkX \verb|max_weight_matching| to the QEC problem
\end{itemize}



\end{document}
